% Shared preamble for main.tex and preview.sh.
\documentclass[10pt,a4paper]{article}
\usepackage[utf8]{inputenc}
\usepackage[T1]{fontenc}
\usepackage{textcomp}
\usepackage{amsmath}
\usepackage{mathpazo}
\usepackage[left=16mm,right=16mm,top=20mm,bottom=22mm]{geometry}
\usepackage{paracol}
\usepackage{microtype}
\usepackage{ragged2e}
\usepackage{xcolor}
\usepackage[hidelinks]{hyperref}
\usepackage{xurl}  % URLs may break anywhere, so long ones stay in the column

\columnratio{0.60}
\setlength{\columnsep}{8mm}
\setlength{\emergencystretch}{3em}  % narrow columns: avoid lines sticking out
\setlength{\parindent}{0pt}
\setlength{\parskip}{0.7\baselineskip}
\definecolor{notecolor}{gray}{0.25}

% ---- Commentary notes -------------------------------------------------------
\newcounter{cnote}
\newcommand{\pendingnotes}{}
% \cnote{kind}{text}: mark the text with a number and queue the note.
\makeatletter
% Two notes in a row get a comma between their marks: each mark ends with a tiny
% kern that the next \cnote can detect.
\newdimen\cnote@kern \cnote@kern=0.01pt
\newcommand{\cnote}[2]{%
  \ifdim\lastkern=\cnote@kern\textsuperscript{\footnotesize,}\fi
  \refstepcounter{cnote}\textsuperscript{\footnotesize\thecnote}\kern\cnote@kern
  %
  \edef\cnote@head{\noexpand\printnote{\thecnote}}%
  \expandafter\g@addto@macro\expandafter\pendingnotes\expandafter{\cnote@head{#1}{#2}}}
\makeatother
\newcommand{\printnote}[3]{%
  {\footnotesize\RaggedRight\color{notecolor}\textsuperscript{#1}\textsc{#2}.\ #3\par}}
% Print queued notes beside the paragraph just finished, then realign the columns.
\newcommand{\flushnotes}{%
  \switchcolumn\pendingnotes\gdef\pendingnotes{}\switchcolumn*}

\newcommand{\cstyle}[1]{\cnote{style}{#1}}
\newcommand{\clogic}[1]{\cnote{logic}{#1}}
\newcommand{\cfact}[1]{\cnote{fact}{#1}}
\newcommand{\ccut}[1]{\cnote{cut}{#1}}
\newcommand{\cauthor}[1]{\cnote{author's footnote}{#1}}
% Paragraph-level note: placed at the start of a paragraph, so it prints first in the margin.
\newcommand{\cpara}[1]{\cnote{paragraph}{#1}}


% Inside the two columns, links print as plain text. A hyperlink that breaks across a
% page boundary inside paracol crashes pdflatex (seen on "On Caring": "\pdfendlink ended
% up in different nesting level", then a segfault). The link to each original, in the
% post header, stays clickable.
\makeatletter
\newcommand{\nolinks}{\let\href\@secondoftwo\let\url\nolinkurl}
\makeatother

% ---- Post header ------------------------------------------------------------
% \post{title}{author}{year}{url} starts a post; \closepost closes it.
\newcommand{\post}[4]{%
  \clearpage
  \phantomsection\addcontentsline{toc}{subsection}{#1}%
  {\noindent\Large #1\par}\vspace{0.3ex}
  {\noindent\small #2, #3.\quad \href{#4}{Original on LessWrong}\par}\vspace{1ex}
  \begin{paracol}{2}\nolinks}
\newcommand{\closepost}{\flushnotes\end{paracol}}

% ---- Summary and Response after each post -------------------------------------
% Single column, same width as the text column, so line length stays readable.
\newenvironment{afterword}
  {\bigskip\begin{list}{}{\leftmargin=0pt\rightmargin=0.4\textwidth
     \listparindent=0pt\itemindent=0pt\parsep=0.7\baselineskip}\item\relax}
  {\end{list}}
\newcommand{\summaryhead}{{\large Summary}\par}
\newcommand{\responsehead}{\medskip{\large Response}\par}

% ---- Elements produced by src/md2tex.py -------------------------------------
% Images are not reproduced; a placeholder links to the original image.
\newcommand{\figph}[1]{{\small\color{notecolor}[Figure in the original: \href{#1}{link}.]}}
% Footnote numbers that the original prints as plain [1], and numbers of end notes.
\newcommand{\fnnum}[1]{\textsuperscript{#1}}
% Section headings inside a post.
\newcommand{\posthead}[1]{\medskip{\large #1}\par}

% Characters pdflatex does not know by default, as they occur in the originals.
\DeclareUnicodeCharacter{00AC}{\ensuremath{\neg}}
\DeclareUnicodeCharacter{00D7}{\ensuremath{\times}}
\DeclareUnicodeCharacter{2248}{\ensuremath{\approx}}
\DeclareUnicodeCharacter{2212}{\ensuremath{-}}
\DeclareUnicodeCharacter{03A0}{\ensuremath{\Pi}}
\DeclareUnicodeCharacter{2260}{\ensuremath{\neq}}
\DeclareUnicodeCharacter{FFFC}{}
\DeclareUnicodeCharacter{0391}{A}
\DeclareUnicodeCharacter{0392}{B}
\DeclareUnicodeCharacter{0394}{\ensuremath{\Delta}}
\DeclareUnicodeCharacter{03A6}{\ensuremath{\Phi}}
\DeclareUnicodeCharacter{21D2}{\ensuremath{\Rightarrow}}
\DeclareUnicodeCharacter{2190}{\ensuremath{\leftarrow}}
\DeclareUnicodeCharacter{2715}{\ensuremath{\times}}
\DeclareUnicodeCharacter{2265}{\ensuremath{\geq}}
% A stray private-use glyph inside a formula in the Bayes essay; it carries no meaning.
\DeclareUnicodeCharacter{10FC06}{}
\DeclareUnicodeCharacter{2215}{/}
\DeclareUnicodeCharacter{2661}{\ensuremath{\heartsuit}}
% Superscript digits used as footnote markers in some originals (¹²³ are in Latin-1 already).
\DeclareUnicodeCharacter{2070}{\textsuperscript{0}}
\DeclareUnicodeCharacter{2074}{\textsuperscript{4}}
\DeclareUnicodeCharacter{2075}{\textsuperscript{5}}
\DeclareUnicodeCharacter{2076}{\textsuperscript{6}}
\DeclareUnicodeCharacter{2077}{\textsuperscript{7}}
\DeclareUnicodeCharacter{2078}{\textsuperscript{8}}
\DeclareUnicodeCharacter{2079}{\textsuperscript{9}}
\DeclareUnicodeCharacter{2208}{\ensuremath{\in}}
\DeclareUnicodeCharacter{221E}{\ensuremath{\infty}}
\DeclareUnicodeCharacter{03A3}{\ensuremath{\Sigma}}
\DeclareUnicodeCharacter{2192}{\ensuremath{\rightarrow}}
\DeclareUnicodeCharacter{2264}{\ensuremath{\leq}}
\DeclareUnicodeCharacter{2009}{\,}
\DeclareUnicodeCharacter{221A}{\ensuremath{\surd}}

% Book and sequence headings (Rationality: A-Z has six books of sequences), with contents
% entries: book = part, sequence = section, post = subsection.
\newcommand{\bookhead}[1]{%
  \clearpage\phantomsection\addcontentsline{toc}{part}{#1}%
  \vspace*{0.3\textheight}{\noindent\huge #1\par}}
\newcommand{\sequencehead}[1]{%
  \clearpage\phantomsection\addcontentsline{toc}{section}{#1}%
  \vspace*{0.2\textheight}{\noindent\LARGE #1\par}}
